\section{A general theorem for elementary Euclidean geometry}
\label{sec:elimination}
The preceding field is more than a convenient place to calculate.
It supports every fixed first-order statement of elementary real geometry.
This section establishes the precise scope of that claim. Its finite
sign-determination argument is closely related to the earlier planar
packet and to formal real-algebraic elimination \cite{planar,cm,pr}.
The theorem is a schema indexed externally by a fixed formula, not an
internal truth predicate for arithmetic.

\subsection{A finite sign calculation}
\begin{lemma}[Signed remainder query]\label{qe:query}
Let $D$ be a squarefree nonconstant polynomial over $\K$ and $W$ any
polynomial. The integer
\[
 Q_D(W)=\sum_{D(r)=0}\sgn W(r)
\]
can be computed using only coefficient arithmetic and sign tests, with
degree bounds on every division and loop.
\end{lemma}
\begin{proof}
Reduce $D'W$ modulo $D$. A zero remainder gives zero query. Otherwise
cancel the gcd of the remainder and $D$, leaving a proper reduced fraction
$B/A$. Form the signed Euclidean remainder chain
\[
 S_0=A,\quad S_1=B,\quad S_{i+1}=-\operatorname{rem}(S_{i-1},S_i).
\]
Stop before the zero remainder. Coprimality makes the final entry a
nonzero constant. Let $V_+$ be the sign-variation count in the leading
coefficients and $V_-$ the count after multiplying each leading sign by
$(-1)^{\deg S_i}$. Then $Q_D(W)=V_--V_+$.

Here is the proof of the formula. List all roots of the product of the
nonzero chain entries using Theorem~\ref{found:roots}. On each intervening
interval every sign is fixed: a contrary sign would produce a root by
polynomial intermediate value. Adjacent entries have no common root,
by the Euclidean identities. At a zero of an interior entry, the
neighboring entries have opposite values because the remainder has a
minus sign. The two adjacent edges therefore contribute exactly one
variation on either side. Simultaneous zeros of nonadjacent entries affect
disjoint pairs of edges. There is no zero in the final constant entry.

At a zero $r$ of $A$, factor $A=(T-r)A_1$, so $A_1(r)=A'(r)$.
The first edge loses one variation if $B(r)/A'(r)>0$ and gains one
if it is negative. The local coefficient of $(T-r)^{-1}$ in $D'W/D$
is $W(r)$; cancellation and removal of a polynomial part preserve that
coefficient. Cancelled roots have $W(r)=0$. Thus the net loss of
variations is the stated sum. Signs at the two tails are their leading
term signs, since $|T|>1+\sum_{i<d}|a_i/a_d|$ suffices for dominance.

Every use of a local sign is quantitative. Expand
$g(r+h)-g(r)=\sum_{j\ge1}c_jh^j$. If $g(r)\ne0$, then
$|h|<\min(1,|g(r)|/(1+\sum|c_j|))$ preserves its sign. Choose a
rational bound below this positive algebraic number using its sign gap.
Consequently the root-list telescoping argument is a bounded Boolean
and integer calculation, with the candidate point kept as a free parameter.
\end{proof}

\begin{lemma}[Every occurring sign pattern is decidable]\label{qe:patterns}
Given polynomials $f_1,\ldots,f_m$ and a squarefree $D$, occurrence of
any prescribed sign vector at a root of $D$ is a PR coefficient test.
Here $D$ is nonzero; a nonzero constant has no roots and every pattern
count is zero.
\end{lemma}
\begin{proof}
On $z\in\{-1,0,1\}$ the polynomials
\[
 I_0(z)=2(1-z^2),\quad I_+(z)=z^2+z,\quad I_-(z)=z^2-z
\]
are twice the indicator of the indicated sign. Expand
$\prod_j I_{\sigma_j}(Z_j)=\sum_e c_eZ^e$. By multiplicativity of
sign and finite distributivity,
\[
 2^m\#\{r:D(r)=0,\ \sgn f_j(r)=\sigma_j\ (1\le j\le m)\}
   =\sum_e c_e Q_D\left(\prod_j f_j^{e_j}\right).
\]
The right side is a computed integer; positivity decides occurrence.
No comparison of unknown roots is needed in this coefficient algorithm.
\end{proof}

\begin{theorem}[One-variable elimination and witnesses]\label{qe:one}
For a quantifier-free ordered-ring formula $\theta(T,\bar a)$, its
existential and universal quantifiers in $T$ can be eliminated by a PR
finite decision tree on polynomial signs in $\bar a$. On an accepting
existential branch, a scalar witness can be constructed. On a rejecting
universal branch, a counterexample can be constructed.
\end{theorem}
\begin{proof}
First fix the parameters. Remove constant polynomials after recording
their signs, including sign zero for identically zero polynomials.
If no nonconstant polynomial remains, $T=0$ is the only
sample required. Otherwise let $P$ be the monic squarefree part of their
product. All possible sign vectors occur in one of three places: at a
root of $P$, at a root of the squarefree part of $P'$, or in either tail.
If $P'$ is a nonzero constant, its root family is empty; the
nonconstant-root query is not called for it.
To justify the middle class, let $a<b$ be consecutive roots of $P$.
Simplicity and the common interior sign of $P$ imply that $P'(a)$ and
$P'(b)$ have opposite signs. Factoring out $T-a$ and $T-b$ makes this
assertion explicit, with the local bounds of Lemma~\ref{qe:query}.
Intermediate value supplies a derivative root in $(a,b)$, where all the
original signs are constant. Hence every bounded open cell is sampled.

The queries in Lemma~\ref{qe:patterns} determine precisely which sign
vectors occur at the two kinds of roots, and leading terms determine the
tails. Evaluate $\theta$ on this finite collection. For actual witnesses
use the sorted roots of $P$ itself, their intervening midpoints, and one
point outside each extreme root. These samples contain every sign vector
needed. If $P$ has no real roots, use the single sample $T=0$ instead:
all original signs are constant by intermediate value, so this sample
is complete even when $P$ is nonconstant.
A true existential has a true sample; a false universal has a
false sample. If all samples satisfy the matrix, an arbitrary candidate
lies in a root singleton or one of the sampled intervals, so it satisfies
the matrix too. This last direct argument is what licenses universal
quantification intuitionistically.

Now treat parameters symbolically. At each coefficient sign test make
three branches, recording equality to zero, positivity, or negativity.
Test every leading coefficient before division, truncate when it vanishes,
and discard identically zero specializations. On a realized branch every
coefficient is a rational function $p/q$ with $q\ne0$; its sign is the
sign of $pq$ because $q^2>0$. Thus all branch guards are polynomial
sign conditions on the parameters. Disjoin the accepting guards.
Unrealizable branches need not be recognized. Division by $a/b$ is used
only on a branch guarded by $a\ne0$, so no actual division by zero
can occur.

For bounds, if there are $m$ polynomials of degree at most $d$, set
$s=1+m(d+1)$. Product and squarefree degrees are below $s$; query
products have degree at most $2md$. A remainder chain takes at most its
initial maximum degree plus two divisions, and each division takes at most
that degree plus two coefficient cancellations. There are $3^m$ sign
patterns and at most $3^m$ monomials in an indicator product. Unroll those
bounded loops and branch at most three ways at each test. Repeated
addition and multiplication of finite expression-tree sizes give a PR
bound on the whole tree. This bound may be enormous; efficiency is not
needed for membership in iPRA. Correctness at a specialized branch uses
the numerical finite-list invariants already described.
\end{proof}

\begin{theorem}[Elementary geometry transfer, formula by formula]
\label{qe:transfer}
For every fixed first-order ordered-field formula $\phi$ there is a PR
Boolean test $b_\phi$ such that iPRA proves
\[
 \operatorname{Valid}(\bar a)\ \longrightarrow\
      \bigl(\phi^\#(\bar a)\leftrightarrow b_\phi(\bar a)=1\bigr).
\]
The standard real field and the coded field $\K$ agree on $\phi$ at
algebraic parameters. Consequently each true closed elementary
Euclidean-geometric sentence has a proof of its numerical translation in
iPRA. This is a metatheorem schema, not a uniform internal assertion of
the soundness of iPRA.
\end{theorem}
\begin{proof}
Eliminate the innermost scalar variable and repeat. The recursion is on
the externally fixed formula. Boolean cases for conjunction, disjunction,
and implication are valid because the relevant subformulas have already
been supplied with decisions. The witness and counterexample arguments
of Theorem~\ref{qe:one} handle both quantifiers directly. Translating
$\forall T\,\theta$ merely as $\neg\exists T\,\neg\theta$ would
be insufficient without that argument; we do not make that omission.

Translate each point binder into two scalar binders, point equality into
component equality, and equal distances into equality of squared distances.
For scalar binders the precise relativization is
\[
 (\exists x\,\psi)^\#=\exists n(\operatorname{Valid}(n)\land\psi^\#),
 \qquad
 (\forall x\,\psi)^\#=\forall n(\operatorname{Valid}(n)\to\psi^\#).
\]
Scalar equality says that the PR comparison test returns zero, not that
the two record numbers are literally equal. All atomic relations respect
this equality by Theorem~\ref{found:field}.
Inclusive betweenness has the polynomial translation
\begin{equation}\label{eq:betweenpoly}
 \Bet(A,X,C)\quad\longleftrightarrow\quad
 [C-A,X-A]=0\ \land\ (X-A)\mathbin\cdot(X-C)\le0.
\end{equation}
Section~\ref{sec:affine} proves this even when $A=C$. Lines and circles
are finite parameter tuples with their validity conditions. Nonnegative
lengths are scalar witnesses satisfying $l^2=N(v)$; uniqueness makes
their introduction a definitional abbreviation. Thus all named elementary
statements in the geometry chapters fall under the translation.

Externally the same elimination tree works over $\R$: polynomial roots
are finite in number, the same sign-variation proof applies, and all
coefficient signs agree on algebraic parameters. At a closed sentence
the final calculation is a closed Boolean computation. Its true value
has a derivation by the defining equations of PR functions, and the
proved equivalence supplies the geometric translation. This proves the
last assertion without introducing a truth predicate into arithmetic.
\end{proof}

\paragraph{A visible example.}
Two circles of squared radius $s$ with centers $(0,0),(1,0)$ intersect
if and only if $4s\ge1$. Subtracting their equations gives $2x=1$;
then $y^2=s-1/4$. The positive, zero, and negative branches give two,
one, and no points, respectively. This is an entire elimination and
witness calculation that can be checked by hand.

\paragraph{What this metatheorem adds.}
It gives a principled route for every further fixed semialgebraic theorem
in the article's neighborhood, rather than an arbitrary cutoff at this
report's named portfolio. It does not assert that every sentence about
continuous curves, arbitrary subsets, integer-indexed geometry, or
transcendental functions is an elementary ordered-field sentence. The
polygon and approximation families below receive their own bounded-loop
arguments when their input length or precision varies.
